\documentclass{article}
\usepackage[margin=1in, a4paper]{geometry}
\usepackage{amsmath, amssymb, amsfonts}
\usepackage{graphicx}
\usepackage{xcolor}
\usepackage{listings}
\usepackage{booktabs}
\usepackage[utf8]{inputenc}
\usepackage[T1]{fontenc}
\usepackage{hyperref}
\hypersetup{colorlinks=true, urlcolor=blue, linkcolor=blue}
\usepackage{enumitem}
\usepackage{float}
\usepackage{longtable}
\usepackage{array}

% Professional color palette for academic publishing
\definecolor{myblue}{cmyk}{0.8,0.4,0,0.1}
\definecolor{mygreen}{cmyk}{0.7,0,0.5,0.2}
\definecolor{myorange}{cmyk}{0,0.5,0.8,0.1}
\definecolor{mygray}{cmyk}{0,0,0,0.4}
\definecolor{arrowcolor}{cmyk}{0.9,0.5,0,0.3}
\definecolor{nodecolor}{cmyk}{0.6,0.2,0,0.05}
\definecolor{framecolor}{cmyk}{0.4,0.1,0,0.1}
% Professional CMYK color palette
\definecolor{darkblue}{cmyk}{1.0,0.7,0,0.3}
\definecolor{darkred}{cmyk}{0,1.0,1.0,0.3}
\definecolor{darkgreen}{cmyk}{0.8,0,1.0,0.4}

% Python listings style definition
\definecolor{codegreen}{rgb}{0,0.6,0}
\definecolor{codegray}{rgb}{0.5,0.5,0.5}
\definecolor{codepurple}{rgb}{0.58,0,0.82}
\definecolor{backcolour}{rgb}{0.99,0.99,0.99}
\lstdefinestyle{mystyle}{
    backgroundcolor=\color{backcolour},   
    commentstyle=\color{codegreen},
    keywordstyle=\color{magenta},
    numberstyle=\tiny\color{codegray},
    stringstyle=\color{codepurple},
    basicstyle=\ttfamily\footnotesize,
    breakatwhitespace=false,         
    breaklines=true,                 
    captionpos=b,                    
    keepspaces=true,                 
    numbers=left,                    
    numbersep=5pt,                  
    showspaces=false,                
    showstringspaces=false,
    showtabs=false,                  
    tabsize=2
}
\lstset{style=mystyle}

\title{\textbf{The Logistics \& Supply Chain Problem-Solving Playbook}}
\author{LaTeX-Bot}
\date{\today}

\begin{document}
\maketitle

\section{The VRP with Split Deliveries (SDVRP)}
The standard Vehicle Routing Problem (VRP) assumes that each customer's demand must be satisfied by a single vehicle visit. The VRP with Split Deliveries (SDVRP) relaxes this constraint, allowing a customer's total demand to be split among multiple vehicles. This seemingly small change dramatically increases the complexity of the problem but can lead to significant cost savings by enabling more efficient vehicle packing and routing. Instead of sending a half-empty truck to satisfy the last few items of a large order, those items can be serviced by another vehicle already passing nearby.

\subsection{The Scenario: Last-Mile Urban Distribution}
Consider a large e-commerce company operating a distribution center on the outskirts of a dense city. Every morning, hundreds of orders with varying sizes and destinations must be delivered. The vehicle fleet is heterogeneous, with vans of different capacities. Standard VRP solvers often produce inefficient routes: a large van might be dispatched with only 80\% of its capacity used because the next undelivered order is too large to fit, forcing a suboptimal route. By allowing deliveries to be split, a large order can be partially delivered by the first van, with the remainder handled by a smaller van that has spare capacity and is heading in the same direction. This improves vehicle utilization, reduces the total number of vehicles needed, and lowers overall mileage and fuel consumption.

\subsection{Tier 1: The Pen \& Paper Method (Mathematical Formulation)}
The SDVRP can be formulated as a Mixed-Integer Linear Program (MILP). The key is to introduce variables that not only decide the path of each vehicle but also the quantity of goods delivered at each stop.

\subsubsection{Model Formulation}
\begin{description}
    \item[Sets]
    \begin{itemize}
        \item $V = \{0, ..., n\}$: The set of nodes, where 0 is the depot and $\{1, ..., n\}$ are the customers.
        \item $K$: The set of available vehicles.
    \end{itemize}

    \item[Parameters]
    \begin{itemize}
        \item $d_i$: The total demand of customer $i \in V \setminus \{0\}$.
        \item $C_k$: The capacity of vehicle $k \in K$.
        \item $c_{ij}$: The travel cost (e.g., distance or time) between node $i$ and node $j$.
    \end{itemize}

    \item[Decision Variables]
    \begin{itemize}
        \item $x_{ijk}$: A binary variable, equal to 1 if vehicle $k$ travels directly from node $i$ to node $j$, and 0 otherwise.
        \item $y_{ik}$: A continuous variable representing the quantity of goods delivered to customer $i$ by vehicle $k$.
    \end{itemize}
\end{description}

\subsubsection{Objective Function}
The objective is to minimize the total travel cost.
\begin{equation}
\text{Minimize} \quad Z = \sum_{i \in V} \sum_{j \in V, i \neq j} \sum_{k \in K} c_{ij} \cdot x_{ijk}
\end{equation}

\subsubsection{Constraints}
\begin{equation}
\sum_{k \in K} y_{ik} = d_i \quad \forall i \in V \setminus \{0\} \label{eq:demand_satisfaction}
\end{equation}

\begin{equation}
\sum_{i \in V \setminus \{0\}} y_{ik} \leq C_k \quad \forall k \in K \label{eq:capacity}
\end{equation}

\begin{equation}
\sum_{j \in V, j \neq i} x_{ijk} = \sum_{j \in V, j \neq i} x_{jik} \quad \forall i \in V, \forall k \in K \label{eq:flow_conservation}
\end{equation}

\begin{equation}
\sum_{j \in V \setminus \{0\}} x_{0jk} \leq 1 \quad \forall k \in K \label{eq:leave_depot}
\end{equation}

\begin{equation}
M \sum_{j \in V, j \neq i} x_{ijk} \geq y_{ik} \quad \forall i \in V \setminus \{0\}, \forall k \in K \label{eq:link_xy}
\end{equation}

\begin{equation}
u_{ik} - u_{jk} + n \cdot x_{ijk} \leq n - 1 \quad \forall i, j \in V \setminus \{0\}, i \neq j, \forall k \in K \label{eq:mtz}
\end{equation}

\begin{equation}
x_{ijk} \in \{0, 1\} \quad \forall i, j \in V, \forall k \in K
\end{equation}

\begin{equation}
y_{ik} \geq 0 \quad \forall i \in V \setminus \{0\}, \forall k \in K
\end{equation}

\begin{equation}
u_{ik} \geq 0 \quad \forall i \in V \setminus \{0\}, \forall k \in K
\end{equation}

Constraint \eqref{eq:demand_satisfaction} ensures each customer's total demand is met by the sum of all deliveries. \eqref{eq:capacity} respects vehicle capacity. \eqref{eq:flow_conservation} ensures that if a vehicle enters a node, it must also leave it. \eqref{eq:leave_depot} ensures each vehicle is used at most once. \eqref{eq:link_xy} is a critical linking constraint: if vehicle $k$ delivers a positive quantity $y_{ik}$ to customer $i$, it must visit customer $i$ (where M is a large number). Finally, \eqref{eq:mtz} represents one form of subtour elimination using auxiliary variables $u_{ik}$.

\begin{figure}[htbp]
\centering
\includegraphics[width=\textwidth]{../figures-png/line76.fig1.png}
\caption{A conceptual SDVRP solution where Customer 3's demand is split.}
\end{figure}

\subsection{Tier 2: The Classic Heuristic (Sweep Algorithm)}
The Sweep Algorithm is a geometrically inspired heuristic well-suited for routing problems. It sequences customers based on their polar angle relative to the depot, creating clusters that can then be formed into routes.

\subsubsection{Algorithm Logic}
\begin{enumerate}
    \item \textbf{Initialization}: Place the depot at the origin of a polar coordinate system. Convert the Cartesian coordinates $(x_i, y_i)$ of each customer $i$ into polar coordinates $(r_i, \theta_i)$.
    \item \textbf{Sweeping}: Start a ray (vector) at an angle of 0. Rotate the ray counter-clockwise. As the ray ``sweeps'' over customers, add them to a list in the order they are encountered.
    \item \textbf{Clustering and Routing}:
    \begin{enumerate}
        \item Iterate through the ordered list of customers.
        \item Start a new route with the first unassigned customer.
        \item Add subsequent customers from the list to the current route as long as vehicle capacity is not exceeded.
        \item \textbf{Split Delivery Logic}: When the next customer's full demand ($d_i$) does not fit, but the vehicle has remaining capacity ($C_{rem} > 0$), create a partial delivery. Assign quantity $C_{rem}$ to the current vehicle for customer $i$, and reduce the customer's remaining demand ($d_i' = d_i - C_{rem}$).
        \item Close the current route and start a new one with the next customer (or the partially-serviced customer if it still has demand).
    \end{enumerate}
    \item \textbf{Improvement}: Once all initial routes are formed, apply a local search heuristic (like 2-opt) to each route to improve its efficiency.
\end{enumerate}

\begin{figure}[htbp]
\centering
\includegraphics[width=\textwidth]{../figures-png/line76.fig2.png}
\caption{The Sweep Algorithm conceptually orders customers by angle to build routes.}
\end{figure}

\subsection{Tier 3: The Advanced Algorithm (Ant Colony Optimization)}
Ant Colony Optimization (ACO) is a metaheuristic inspired by the foraging behavior of ants. It is particularly effective for graph-based optimization problems like the VRP. For SDVRP, the ant's decision-making process must be enhanced.

\subsubsection{ACO for SDVRP}
In ACO, artificial ants construct solutions (routes) by probabilistically moving between nodes. The path choice is influenced by two factors:
\begin{itemize}
    \item \textbf{Pheromone ($\tau_{ij}$)}: The accumulated "memory" of good solutions on the arc between node $i$ and $j$.
    \item \textbf{Heuristic Information ($\eta_{ij}$)}: A greedy measure, typically the inverse of the cost ($1/c_{ij}$).
\end{itemize}

An ant $k$ at node $i$ chooses to go to node $j$ with probability:
\begin{equation}
p_{ij}^k = \frac{[\tau_{ij}]^\alpha [\eta_{ij}]^\beta}{\sum_{l \in \text{Allowed}_k} [\tau_{il}]^\alpha [\eta_{il}]^\beta}
\end{equation}
where $\alpha$ and $\beta$ control the importance of pheromone vs. heuristic information.

\subsubsection{Handling Split Deliveries}
The core modification for SDVRP lies in the ant's state and decision process. When an ant arrives at a customer node, it must decide how much to deliver.
\begin{enumerate}
    \item \textbf{Ant State}: Each ant maintains its current location, remaining capacity, and the list of customers already visited in its current tour.
    \item \textbf{Decision at a Customer Node}: When an ant visits customer $i$, it delivers the maximum possible amount: $\min(\text{remaining\_demand}_i, \text{ant\_capacity}_{rem})$.
    \item \textbf{Modifying the Node List}: After the delivery, if customer $i$ still has remaining demand, it is NOT removed from the list of `Allowed` nodes. It can be visited again by the same or a different ant. This is the key difference from standard VRP ACO.
    \item \textbf{Pheromone Update}: After all ants complete their tours, the pheromone trails are updated. Trails belonging to shorter total routes (better solutions) receive more pheromone. A global update rule is:
    \begin{equation}
    \tau_{ij} \leftarrow (1-\rho)\tau_{ij} + \sum_{k} \Delta \tau_{ij}^k
    \end{equation}
    where $\rho$ is the evaporation rate and $\Delta \tau_{ij}^k$ is the pheromone deposited by ant $k$, usually inversely proportional to its tour length.
\end{enumerate}

\begin{figure}[htbp]
\centering
\includegraphics[width=\textwidth]{../figures-png/line76.fig3.png}
\caption{ACO for SDVRP. Node C3 is visited by two different ants to satisfy its demand.}
\end{figure}

\subsection{Tier 4: The AI/ML/RL Augmentation Method}
For a problem as dynamic as SDVRP, Reinforcement Learning (RL) combined with Graph Neural Networks (GNNs) offers a powerful, adaptive approach. The model learns a routing policy directly from experience, making it robust to changing demand patterns and operational conditions.

\subsubsection{RL Formulation for SDVRP}
\begin{itemize}
    \item \textbf{State (S)}: The state is a graph representation of the current situation. It includes:
    \begin{itemize}
        \item Node features: Remaining demand for each customer, location coordinates.
        \item Vehicle features: Current location, remaining capacity.
        \item Global features: Time of day, overall progress.
    \end{itemize}
    \item \textbf{Action (A)}: At each decision step, the RL agent (representing a vehicle) chooses an action. A composite action is required:
    \begin{itemize}
        \item Action Part 1: Select the next customer to visit from the set of uncompleted customers.
        \item Action Part 2: Decide the quantity to deliver to that customer (from a discrete set of options, e.g., \{25\%, 50\%, 75\%, 100\% of remaining need\} or up to remaining capacity).
    \end{itemize}
    \item \textbf{Reward (R)}: The reward function guides the learning process. The agent receives a large negative reward proportional to the total distance traveled at the end of an episode (when all deliveries are complete). Small positive rewards could be given for completing a delivery, incentivizing progress.
    \begin{equation}
    R = -(\text{total distance}) - \lambda \cdot (\text{number of vehicles used})
    \end{equation}
\end{itemize}

\subsubsection{Graph Neural Network Architecture}
A GNN is used as the policy network's function approximator. It naturally processes the graph structure of the VRP.
\begin{enumerate}
    \item \textbf{Graph Embedding}: A GNN takes the state graph as input. Through layers of message passing, it learns a rich embedding (a feature vector) for each node (depot, customers, vehicles) that captures its context within the current routing plan.
    \item \textbf{Policy Head}: The node embeddings are fed into a policy head (e.g., an attention mechanism). This head computes a probability distribution over the possible next actions. For SDVRP, it would output probabilities for each `(customer, quantity)` pair.
    \item \textbf{Training}: The model is trained using a policy gradient algorithm like PPO (Proximal Policy Optimization) or A2C (Advantage Actor-Critic). The agent explores by taking actions, receives rewards, and updates the GNN's weights to increase the probability of taking actions that lead to higher cumulative rewards (lower costs).
\end{enumerate}

\begin{figure}[htbp]
\centering
\includegraphics[width=\textwidth]{../figures-png/line76.fig4.png}
\caption{A GNN-based Reinforcement Learning architecture for SDVRP.}
\end{figure}

\subsection{Tier 5: The Integrated Digital Twin}
A Digital Twin for SDVRP is a high-fidelity, real-time virtual replica of the entire delivery ecosystem. It moves beyond static planning to dynamic, continuous optimization.

The twin integrates:
\begin{itemize}
    \item \textbf{Static Data}: Road networks, customer locations, vehicle specifications, distribution center layouts.
    \item \textbf{Real-Time Data Streams}: GPS data from vehicles, live traffic feeds (e.g., from Google Maps APIs), weather updates, new incoming orders, and order cancellations.
\end{itemize}

\textbf{Capabilities}:
\begin{enumerate}
    \item \textbf{Live Re-optimization}: The initial plan is just a baseline. If a vehicle gets stuck in traffic, its ETA for subsequent stops grows. The twin detects this deviation and automatically triggers a re-optimization. It might decide to split a future delivery from the delayed truck's route and assign it to a nearby, on-schedule truck to maintain service level agreements (SLAs).
    \item \textbf{Predictive What-If Analysis}: Planners can inject hypothetical events into the twin to test resilience. ``What if the main bridge closes for 2 hours this afternoon?'' The twin simulates the cascading delays and computes the optimal contingency plan, including dynamic re-splitting of deliveries across the fleet.
    \item \textbf{Capacity Planning}: By simulating historical data and future demand forecasts, the twin can accurately predict fleet requirements, identifying the need for more vehicles or highlighting opportunities to reduce the fleet size.
\end{enumerate}

\subsection{Tier 6: The Autonomous \& Self-Optimizing Ecosystem}
This tier envisions a decentralized system where vehicles are autonomous agents that negotiate and collaborate to fulfill deliveries. The central planner's role shifts from command-and-control to setting system-wide goals.

\textbf{Multi-Agent System (MAS) for SDVRP}:
\begin{enumerate}
    \item \textbf{Agents}: Each vehicle is an intelligent agent. The depot or a central server acts as an "auctioneer".
    \item \textbf{Task Broadcasting}: The depot broadcasts a list of outstanding delivery tasks, where each task is a (customer, quantity) pair. A large order for customer C1 with demand 100 might be broadcast as ten tasks of (C1, 10 units).
    \item \textbf{Bidding Mechanism}: Each vehicle agent calculates a "bid" for each task. The bid reflects the marginal cost for the agent to add that task to its current planned route. This cost is a function of extra distance, time, and fuel.
        \begin{equation}
        \text{Bid}_{ik} = \text{Cost}(\text{Route}_k \cup \text{Task}_i) - \text{Cost}(\text{Route}_k)
        \end{equation}
    \item \textbf{Winner Determination}: The auctioneer agent awards each task to the lowest bidder. A single customer's demand can be won by multiple vehicle agents, naturally resulting in a split delivery.
    \item \textbf{Continuous Re-bidding}: The system is not static. As vehicles move and conditions change, agents can drop tasks (with a penalty) and re-bid on others, leading to continuous, emergent optimization of the entire delivery network. This decentralized approach is extremely robust to disruption.
\end{enumerate}

\subsection{Tier 7: The Human-AI Symbiotic Partnership (The Centaur Model)}
In the Centaur model, the AI is not a replacement for the human planner but an augmentation tool, creating a powerful human-AI team. The focus is on leveraging the AI's computational power and the human's intuition and domain knowledge.

An interactive dashboard is the centerpiece of this collaboration:
\begin{itemize}
    \item \textbf{AI-Generated Baseline}: The system (e.g., the Tier 4 RL model) generates an optimal SDVRP plan. This is presented visually on a map.
    \item \textbf{Interactive Manipulation}: The human planner can override the AI. They can drag a stop from one vehicle's route to another, manually change the split quantity for a customer, or re-sequence stops within a route.
    \item \textbf{Real-Time Impact Analysis}: As the planner makes a change, the AI instantly re-calculates and displays the impact. A sidebar might update with: ``Total Distance: +8.5km, ETA for Customer X: -15 min, Vehicle 3 Capacity: 98\% -> 112\% (Overloaded!)''.
    \item \textbf{Explainable AI (XAI)}: The most crucial component. The planner can query the AI's decisions.
        \begin{itemize}
            \item \textbf{Planner}: ``Why was Customer Y's delivery split?''
            \item \textbf{AI}: ``Vehicle 1 serves 80\% of Customer Y's order as it passes within 0.5km. Assigning the final 20\% to Vehicle 4, which is also nearby, avoids a 12km detour for Vehicle 1, saving 25 minutes and \$5 in fuel costs.''
        \end{itemize}
\end{itemize}
This symbiotic loop allows for fine-tuning based on unquantifiable factors (e.g., a specific customer's preferences, known roadwork not in any live feed) while still benefiting from algorithmic optimization.

\subsection{Tier 8: The Value-Aligned \& Ethical Framework}
An optimization algorithm, if purely focused on minimizing cost, can create solutions that are efficient but unfair or undesirable in other ways. An ethical framework embeds values directly into the optimization process.

For SDVRP, key ethical considerations include:
\begin{enumerate}
    \item \textbf{Driver Fairness}: A pure cost-minimizing solution might consistently give one driver all the difficult, high-traffic, long-distance routes while another gets easy, clustered, suburban routes. This can lead to disparities in workload, stress, and (if paid per delivery) income.
        \begin{itemize}
            \item \textbf{Ethical Constraint}: Add constraints to the model (Tier 1) or penalty terms to the objective function to balance workloads. For example: 
            \begin{equation}
            \max_k(\text{TourLength}_k) - \min_k(\text{TourLength}_k) \leq \Delta_{max}
            \end{equation}
            This ensures the difference in tour length between any two drivers is bounded.
        \end{itemize}
    \item \textbf{Customer Equity}: The algorithm might de-prioritize deliveries to low-profitability or remote areas, leading to consistently later delivery windows for certain communities. Split deliveries could exacerbate this if remote customers always get the "leftover" part of a delivery at the end of the day.
        \begin{itemize}
            \item \textbf{Ethical Objective}: Introduce a multi-objective function that balances cost with a fairness metric, such as minimizing the variance of delivery times across all customers.
        \end{itemize}
    \item \textbf{Environmental Impact}: Minimizing distance is a good proxy for minimizing fuel, but not perfect. An ethical framework might add a specific objective term to minimize total carbon emissions, which could favor electric vehicles for certain routes or penalize routes with significant idling time in traffic, even if they are shorter.
\end{enumerate}
The system's "constitution" would define these trade-offs, ensuring that the pursuit of efficiency doesn't violate core principles of fairness and sustainability.

\subsection{Tier 9: The Quantum Leap (Computational Supremacy)}
Combinatorial optimization problems like SDVRP are NP-hard, meaning their complexity explodes as the number of customers and vehicles grows. Quantum computing offers a fundamentally new paradigm for tackling such problems. The Quantum Approximate Optimization Algorithm (QAOA) is a leading candidate for solving VRP-like problems on near-term quantum hardware.

\subsubsection{SDVRP as a QUBO/QAOA Problem}
\begin{enumerate}
    \item \textbf{Qubit Encoding}: First, the problem's decision variables must be mapped to qubits. The binary routing variable $x_{ijk}$ (vehicle $k$ travels from node $i$ to $j$) is a natural fit. For a problem with $N$ customers and $K$ vehicles, we could have a qubit for each possible edge in each vehicle's tour. Let's simplify and use $x_{ipk}$ which is 1 if customer $i$ is at position $p$ in the tour of vehicle $k$. The split delivery quantities add another layer of complexity, often handled by binning the demand into discrete units.

    \item \textbf{Cost Hamiltonian ($H_C$)}: The objective function is encoded into a "Cost" or "Problem" Hamiltonian. This is a diagonal operator where the eigenvalues correspond to the cost of each possible solution (a specific configuration of qubit states). For SDVRP, $H_C$ would be a weighted sum of terms, where each term corresponds to a travel cost $c_{ij}$ multiplied by the qubit variables representing that travel segment ($x_{i,p,k}$ and $x_{j,p+1,k}$).
    \begin{equation}
    H_C = \sum_{i,j,p,k} c_{ij} \cdot x_{i,p,k} \cdot x_{j,p+1,k} + \text{Penalty Terms}
    \end{equation}
    Penalty terms are added to enforce constraints, such as vehicle capacity and ensuring each customer is visited.

    \item \textbf{Mixer Hamiltonian ($H_M$)}: A "Mixer" Hamiltonian is used to enable transitions between different solution states, allowing the algorithm to explore the solution space. A common choice is the transverse field Hamiltonian, $H_M = \sum_i \sigma_i^X$, where $\sigma_i^X$ is the Pauli-X operator on qubit $i$.

    \item \textbf{QAOA Algorithm}:
    \begin{enumerate}
        \item Start with the qubits in an equal superposition of all possible states (created by applying Hadamard gates).
        \item Repeatedly apply the cost Hamiltonian $e^{-i\gamma_l H_C}$ and the mixer Hamiltonian $e^{-i\beta_l H_M}$ for a chosen number of layers, $p$.
        \item The angles $(\gamma_1, ..., \gamma_p)$ and $(\beta_1, ..., \beta_p)$ are classical parameters optimized by an outer classical loop (e.g., a gradient-free optimizer). The goal is to find angles that maximize the probability of measuring a low-cost state.
        \item Finally, measure the state of the qubits. The resulting bitstring represents a candidate solution to the SDVRP.
    \end{enumerate}
\end{enumerate}

\begin{figure}[htbp]
\centering
\includegraphics[width=\textwidth]{../figures-png/line76.fig5.png}
\caption{Conceptual structure of a QAOA circuit for solving an optimization problem. The angles $\gamma$ and $\beta$ are optimized classically to guide the quantum state toward a low-cost solution.}
\end{figure}

\newpage
\subsection{Practice Questions}

\begin{enumerate}
    \item \textbf{(Tier 1)} You have two vehicles, each with a capacity of 100 units. You need to deliver to two customers: Customer A with a demand of 70, and Customer B with a demand of 130. Set up the demand satisfaction constraint (Eq. \ref{eq:demand_satisfaction}) and the linking constraint (Eq. \ref{eq:link_xy}) for Customer B.
    
    \item \textbf{(Tier 2)} Given a set of customers arranged in a perfect circle around a depot, describe how the Sweep Algorithm would cluster them. What is the main drawback of this method if the customers' demands are highly variable?
    
    \item \textbf{(Tier 3)} In the ACO model for SDVRP, what would be the effect of setting the pheromone evaporation rate $\rho$ to a very high value (e.g., 0.9)? What about a very low value (e.g., 0.01)?
    
    \item \textbf{(Tier 4)} For the GNN-based RL agent, if a vehicle is half-full and located near two customers (A and B), where A has a large remaining demand and B has a small one, how would the node embeddings for A and B likely differ? How would that influence the policy network's output?
    
    \item \textbf{(Tier 5)} Propose a "what-if" scenario for a Digital Twin of a grocery delivery service that would be very difficult to analyze with static planning tools. What data streams would be essential for the twin to evaluate your scenario?
    
    \item \textbf{(Tier 6)} In the Multi-Agent System for SDVRP, imagine two vehicles are equidistant from a customer. Vehicle 1 is 90\% full, and Vehicle 2 is 20\% full. Which vehicle is likely to win the bid to service that customer's demand, and why? Explain your reasoning based on the marginal cost calculation.
    
    \item \textbf{(Tier 7)} You are designing a Centaur dashboard for SDVRP. A logistics planner sees that the AI has split a delivery for a VIP customer. The planner wants to force a single delivery for this customer. What information should the XAI provide immediately after the planner makes this manual override?
    
    \item \textbf{(Tier 8)} An SDVRP algorithm consistently assigns partial, end-of-day deliveries to a specific rural region, causing residents there to always receive their packages late. How could you formulate an ethical constraint or objective to mitigate this systemic bias?
    
    \item \textbf{(Tier 9)} In a QAOA formulation for SDVRP, what physical phenomenon allows the algorithm to explore many possible routes simultaneously, and what is the role of the final measurement step?
\end{enumerate}

\end{document}